\documentclass[../../main.tex]{subfiles}

\begin{document}

\section{Variational fermionic wave functions on lattices}
\label{sec:variational-fermionic-wave-functions}

\subsection{Summary}
% (left blank in the handwritten notes -- fill in once part 2 exists)

\subsection{Crash course on the origin of antisymmetric wave functions}

This can be found in probably any introductory book on quantum mechanics, here I directly follow Schwabl, \emph{Quantum Mechanics}, section 13.1.

The Heisenberg uncertainty relation implies that identical particles are
indistinguishable. Consequently the probability amplitude
$|\psi(x_1,\dots,x_N)|^2$ of a quantum many-body state is invariant under the
exchange $x_i\leftrightarrow x_j$ of two indistinguishable particles. Define the
exchange operator
\begin{equation}
P_{ij}\,\psi(x_1,\dots,x_i,\dots,x_j,\dots,x_N)
\coloneqq \psi(x_1,\dots,x_j,\dots,x_i,\dots,x_N),
\end{equation}
so that invariance of the amplitude gives
\begin{equation}
P_{ij}\,\psi(\vec x)=\lambda\,\psi(\vec x),\qquad |\lambda|^2=1 .
\end{equation}
Furthermore, exchanging the same two particles twice cannot change the wave
function at all, hence
\begin{equation}
P_{ij}^2\,\psi(\vec x)=\lambda^2\psi(\vec x)=\psi(\vec x)
\quad\Longrightarrow\quad \lambda^2=1,\ \text{i.e.}\ \lambda=\pm1 .
\end{equation}
Fermions are the particles for which $\lambda=-1$, i.e.
\begin{equation}
\boxed{\;\psi(\dots,x_i,\dots,x_j,\dots)=-\,\psi(\dots,x_j,\dots,x_i,\dots)\;}
\end{equation}
and thus $\psi(\dots,x_i,\dots,x_i,\dots)=0$.

\subsection{Slater determinants}

Here we loosely follow E.\ Koch, \emph{Mean-Field Theory: Hartree-Fock and BCS}
(J\"ulich 2024), section 1.

Let $\{\varphi_n\}$ be a complete, orthonormal single-particle basis,
\begin{align}
\sum_n \varphi_n^*(x)\varphi_n(x')&=\delta(x-x') &&\text{(complete)},\\
\int\! dx\, \varphi_n^*(x)\varphi_m(x)&=\delta_{nm} &&\text{(orthonormal)}.
\end{align}
An arbitrary $N$-particle function can then be expanded in its first argument as
\begin{align}
a(x_1,\dots,x_N)&=\sum_{n_1} a_{n_1}(x_2,\dots,x_N)\,\varphi_{n_1}(x_1),\qquad\text{with}\\
a_{n_1}(x_2,\dots,x_N)&=\int\! dx'_1\, \varphi^*_{n_1}(x'_1)\,a(x'_1,x_2,\dots,x_N).
\end{align}
Iterating in all arguments, yields
\begin{equation}
\psi(x_1,\dots,x_N)=\sum_{n_1,\dots,n_N} a_{n_1,\dots,n_N}\,
\varphi_{n_1}(x_1)\cdots\varphi_{n_N}(x_N),
\end{equation}
which is not antisymmetric yet.
Now apply the antisymmetrizer $\mathcal A$, with $\mathrm{sgn}(P)=(-1)^P$:
\begin{align}
\mathcal A\,\psi(x_1,\dots,x_N)
&\overset{\text{[K~(3)]}}{\equiv}  \frac{1}{\sqrt{N!}}\sum_P (-1)^P\,\psi(x_{P(1)},\dots,x_{P(N)})\\
&=\frac{1}{\sqrt{N!}}\sum_P (-1)^P \sum_{n_1,\dots,n_N} a_{n_1,\dots,n_N}\,
\varphi_{n_1}(x_{P(1)})\cdots\varphi_{n_N}(x_{P(N)})\\
&=\sum_{n_1,\dots,n_N} a_{n_1,\dots,n_N}\,\Phi_{n_1,\dots,n_N}(x_1,\dots,x_N),
\end{align}
where we have used the \emph{Slater determinant}
\begin{align}
\Phi_{\alpha_1,\dots,\alpha_N}(x_1,\dots,x_N)
&\overset{\text{[K~(4)]}}{\equiv} \frac{1}{\sqrt{N!}}
\begin{vmatrix}
\varphi_{\alpha_1}(x_1) & \varphi_{\alpha_2}(x_1) & \cdots & \varphi_{\alpha_N}(x_1)\\
\varphi_{\alpha_1}(x_2) & \varphi_{\alpha_2}(x_2) & \cdots & \varphi_{\alpha_N}(x_2)\\
\vdots & \vdots & \ddots & \vdots\\
\varphi_{\alpha_1}(x_N) & \varphi_{\alpha_2}(x_N) & \cdots & \varphi_{\alpha_N}(x_N)
\end{vmatrix}\\
&\overset{\text{Leibniz}}{=}\frac{1}{\sqrt{N!}}\sum_{P\in S_N}
\mathrm{sgn}(P)\prod_{i=1}^{N}\varphi_{\alpha_i}\big(x_{P(i)}\big).
\end{align}
Hence \textbf{any $\psi(x_1,\dots,x_N)$ can be expanded in Slater determinants.}
Furthermore, it follows immediately that $\psi(\vec x)$ vanishes if either
$x_i=x_j$ \underline{or} $\alpha_i=\alpha_j$ for $i\neq j$:
\textbf{fermions cannot occupy the same position $x$ or the same quantum state
$\alpha$.}

\subsection{The variational principle}

Following F.\ Becca \& S.\ Sorella, \emph{Quantum Monte Carlo Approaches for
Correlated Systems}, section 1.4.

Formulate some variational state $|\psi_{\mathrm{var}}\rangle$ and consider
\begin{equation}
E_{\mathrm{var}}=\frac{\langle\psi_{\mathrm{var}}|H|\psi_{\mathrm{var}}\rangle}
{\langle\psi_{\mathrm{var}}|\psi_{\mathrm{var}}\rangle}
\qquad\text{[B\&S~(1.38)]}.
\end{equation}
Expanding in the energy eigenstates $|\varphi_i\rangle$,
\begin{equation}
|\psi_{\mathrm{var}}\rangle=\sum_i |\varphi_i\rangle\langle\varphi_i|\psi_{\mathrm{var}}\rangle
=\sum_i a_i|\varphi_i\rangle
\qquad\text{[B\&S~(1.39)]},
\end{equation}
so that, using $\langle\varphi_i|\varphi_j\rangle=\delta_{ij}$,
\begin{align}
E_{\mathrm{var}}&=\Big(\sum_i|a_i|^2\Big)^{-1}
\Big(\sum_i a_i|\varphi_i\rangle\Big)^{\dagger} H
\Big(\sum_j a_j|\varphi_j\rangle\Big)
=\Big(\sum_i|a_i|^2\Big)^{-1}\sum_i |a_i|^2 E_i .
\end{align}
Assume $E_i\le E_j$ for $i<j$, i.e.\ $E_0$ is the ground-state energy. Then
\begin{equation}
E_{\mathrm{var}}-E_0=\Big(\underbrace{\sum_i|a_i|^2}_{>0}\Big)^{-1}
\Big(\sum_{i\neq0}\underbrace{|a_i|^2}_{\ge0}
\underbrace{(E_i-E_0)}_{\ge0}\Big)\ \ge\ 0
\qquad\text{[B\&S~(1.41)]},
\end{equation}
i.e.\ \textbf{$E_{\mathrm{var}}$ is an upper bound for the exact energy.} In
practice: parametrize some $|\psi_{\mathrm{var}}\rangle$ and minimize
$E_{\mathrm{var}}$.

\subsection{Variational wave functions}

\subsubsection{Hartree-Fock wave functions}

The Hartree-Fock (HF) approximation writes the many-body wave function as a
product of single-particle orbitals,
\begin{equation}
\boxed{\;
|\psi_{\mathrm{HF}}\rangle=\prod_{\alpha=1}^{N_e}\phi^\dagger_\alpha|0\rangle,
\qquad
\phi^\dagger_\alpha=\sum_{i=1}^{N}\sum_\sigma W_{\sigma,\alpha,i}\,c^\dagger_{i,\sigma},
\quad \sigma=\uparrow,\downarrow \;}
\end{equation}
with $N$ lattice sites and $N_e$ electrons. This leaves
$2\cdot N_e\cdot N$ parameters instead of the exponentially many amplitudes of a
generic many-body state.
% As written in the handwritten notes: "2 N_e N_e instead of 2^{N_e}"; the point
% is 2 (spin) x N_e (orbitals) x N (sites) vs. exponentially many parameters.
The parameters are then optimized by minimizing
\begin{equation}
E_{\mathrm{HF}}=\frac{\langle\psi_{\mathrm{HF}}|H|\psi_{\mathrm{HF}}\rangle}
{\langle\psi_{\mathrm{HF}}|\psi_{\mathrm{HF}}\rangle},
\qquad\text{e.g.}\quad
W_{\sigma,\alpha,i}\ \to\ W_{\sigma,\alpha,i}
-\frac{\partial E_{\mathrm{HF}}}{\partial W^*_{\sigma,\alpha,i}}\,\delta\tau
\qquad\text{[B\&S~(1.54)]},
\end{equation}
using Wirtinger derivatives (see \cref{sec:wirtinger-derivatives}).

The HF wave function can equivalently be written in first quantization as
\textbf{one Slater determinant}. Since it is a \emph{product of single-particle
orbitals}, it is mostly useful in the non-interacting case or at weak coupling.
In the Hubbard model, turning up $U$ makes the motion of one electron depend on
where the other electrons are, i.e.\ it generates electron-electron correlation.

\subsubsection{The Gutzwiller wave function}

Gutzwiller's idea: describe the effect of the Hubbard $U$ interaction by reducing
the weight of configurations with multiply occupied sites. Start from some ground
state $|\psi_0\rangle$ (e.g.\ non-interacting, a Slater determinant) and apply an
operator $P_g$ that suppresses such configurations:
\begin{align}
|\psi_g\rangle &= P_g|\psi_0\rangle\;  \qquad \text{[B\&S~(1.58)]} \\
\text{with} \qquad P_g &=\exp\Big[-\frac{g}{2}\sum_{i=1}^{N}(n_i-n)^2\Big]
\qquad\text{[B\&S~(1.59)]},
\end{align}
where $n_i=\sum_\sigma n_{i\sigma}$ and
$n=\frac{1}{N}\big\langle\sum_{i=1}^{N} n_i\big\rangle$ is the average density per
site. Here $g$ is the variational parameter ($g>0$ for the repulsive Hubbard
model).

Given some configuration $|x\rangle$, $P_g$ counts how many sites within that
state have higher or lower occupation than the average; for repulsive
interactions these states are suppressed accordingly:
\begin{equation}
\langle x|\psi_g\rangle=P_g(x)\,\langle x|\psi_0\rangle,
\qquad P_g(x)\le 1
\qquad\text{[B\&S~(1.60)]}.
\end{equation}
Since $P_g(x)\ge 0$, it \emph{does not} affect the sign structure.

In the Hubbard model a metal-insulator transition is expected at \emph{finite}
$U/t$ for $n=1$. The Gutzwiller wave function can only reproduce this for
$g\to\infty$; otherwise there is always a non-zero probability for
doublon-holon creation, and these could then move around freely and thus conduct.

\paragraph{Fully projected Gutzwiller wave function ($g\to\infty$).}
Examine $P_\infty$ for $n=1$:
\begin{equation}
P_\infty=\lim_{g\to\infty}P_g
=\lim_{g\to\infty}\exp\Big[-\frac{g}{2}\sum_{i=1}^{N}(\hat n_i-1)^2\Big],
\end{equation}
\begin{align}
\Rightarrow\quad
P_\infty|\psi\rangle &=\lim_{g\to\infty}
\begin{cases}
\exp\!\big(-\tfrac{g}{2}a\big), & \text{if }\exists\ \text{site with a holon/doublon }(a>0) \notag\\[2pt]
\exp(0), & \text{else}
\end{cases}
\ |\psi\rangle \\
&=\begin{cases}0\\ |\psi\rangle\end{cases}
\end{align}
Therefore we can rewrite $P_\infty$ as (B\&S~(1.63))
\begin{equation}
P_\infty=\prod_i\big(n_{i,\uparrow}-n_{i,\downarrow}\big)^2.
\end{equation}
If we want to allow empty sites, one has to set
\begin{equation}
P_\infty=\prod_i\big(1-n_{i,\uparrow}n_{i,\downarrow}\big),
\end{equation}
which is important for $n<1$, since then holons and doublons do not come in
pairs.

For $g=\infty$ no density fluctuations can occur, whereas in a true insulator
density fluctuations do occur but produce no net current -- e.g.\ double-occupancy
fluctuations via virtual processes in the Hubbard model at large $U$\footnote{see \textit{Lecture Notes on Correlation and Magnetism}, Patrik Fazekas, chapter 5, $t$-$J$ model)}.

\subsubsection{The density-density Jastrow wave function}

Since the Gutzwiller wave function does not correctly describe an insulator, the
ansatz requires modification. We can, for instance, change the correlation term
that is applied to $|\psi_0\rangle$ (or change $|\psi_0\rangle\rightsquigarrow
|\psi_{U=\infty}\rangle$). The Jastrow factor does this by including long-range
correlations:
\begin{align}
\mathcal J=\exp\Big[-\frac{1}{2}\sum_{i,j}v_{ij}\,(n_i-n)(n_j-n)\Big]
&\qquad\text{[B\&S~(1.65)]}\\
|\psi_{\mathcal J}\rangle=\mathcal J|\psi_0\rangle=|\psi_g\rangle,
\quad\text{for } v_{ij}=\delta_{ij}\,g
&\qquad\text{[B\&S~(1.64)]}.
\end{align}
For translationally invariant systems the variational parameters are
$v_{ij}=v_{|i-j|}$, i.e.\ $\mathcal O(L)$ possible different distances.

The long-range correlations induced by $\mathcal J$ shall create a bound state
between holons and doublons, possibly hindering conduction while still allowing
local density fluctuations. It can be shown that the behaviour of $v_{ij}$
resulting from an energy minimization is \emph{not} enough to capture a fully
gapped phase (power-law behaviour is present, but the Fourier transform of
$v_{ij}$ decays exponentially).

\paragraph{Jastrow-Slater wave function.}
In the case $v_{ij}=v_{ji}$ the Jastrow factor is symmetric under particle
exchange, so (in first quantization) it can be combined with a Slater determinant
to obtain a suitable fermionic wave function,
$\psi_{\mathrm{JS}}(\vec x)=\langle\vec x|\psi_{\mathrm{JS}}\rangle$:
\begin{equation}
\psi_{\mathrm{JS}}(\vec r_1,\dots,\vec r_{N_e})
=\exp\Big[-\sum_{i<j}v_{ij}(n_i-n)(n_j-n)\Big]\cdot
\psi_{\mathrm{HF}}(\vec r_1,\dots,\vec r_{N_e}),
\end{equation}
with $\psi_{\mathrm{HF}}(\vec r_1,\dots,\vec r_{N_e})=\det\{\phi_\alpha(\vec r_i)\}$
and $\{\phi_\alpha(\vec r_i)\}$ a set of one-particle orbitals (see B\&S (1.68)).

\subsubsection{The backflow wave function}

As explained earlier, the set of Slater determinants forms a basis of the
$N$-electron Hilbert space, so one can systematically improve HF by adding more
Slater determinants, but this is computationally very demanding. Instead:
\textbf{try to improve one single Slater determinant by making the orbitals
configuration dependent}, which goes beyond mean field.

\paragraph{Original idea.}
The movement of one particle ``pushes other particles away'' in order to avoid
large overlap among them, i.e.\ a ``counter flow'' of all other particles --
hence \emph{backflow correlations}. One therefore constructs the Slater
determinant with modified coordinates (continuum backflow),
\begin{equation}
\vec r_i^{\,b}=\vec r_i+\sum_{j\neq i}\eta\big(|\vec r_i-\vec r_j|\big)
\big(\vec r_i-\vec r_j\big),
\end{equation}
which also describes the effective displacement of the $i$-th particle due to the
$j$-th one.

\paragraph{On a lattice.}
For lattice models this effective displacement is much less useful, since all
distances are discretized. Instead one changes the one-particle orbitals:
\begin{equation}
\phi_k\big(\vec r^{\,b}_{i,\sigma}\big)\ \simeq\ \phi^b_k\big(\vec r_{i,\sigma}\big)
\equiv \phi_k\big(\vec r_{i,\sigma}\big)+\sum_j c_{ij}[x]\,\phi_k\big(\vec r_{j,\sigma}\big)
\end{equation}
where $\phi_k(\vec r_{i,\sigma})=\langle 0|c_{i,\sigma}|\phi_k\rangle$, the
$|\phi_k\rangle$ are eigenstates of the mean-field Hamiltonian, and $|x\rangle$ is
the many-body configuration of the system.

This is essentially a \emph{dynamic linear combination of orbitals}: dynamic,
since it \emph{depends on the electron configuration}. In the Hubbard model at
large $U$ one can hybridize the orbitals as follows:
\begin{equation}
\phi^b_k(\vec r_{i,\sigma})=\varepsilon\,\phi_k(\vec r_{i,\sigma})
+\eta\sum_j t_{ij}\,\big(D_i H_j\big)\,\phi_k(\vec r_{j,\sigma})
\end{equation}
i.e.\ hybridize the orbitals at $i$ and $j$ if there is a doublon at $i$ and a
holon at $j$, such that the electron can minimize energy by going to the empty
site\footnote{cf.\ \emph{``Role of
backflow correlations for the non-magnetic phase of the $t$-$t'$ Hubbard model''} by Tocchio and Sorella, Eqs.~(2)-(4)}.

\paragraph{Summary.}
Backflow correlations improve a single Slater determinant by replacing fixed
single-particle orbitals with \textbf{configuration-dependent orbitals},
effectively introducing \textbf{many-body corrections} through a determinant whose
elements depend on the full configuration. Combining this with a Jastrow factor
gives the Slater-Jastrow-backflow wave function.

\subsection{Parametrization via neural networks}

Now let us have a look at how neural networks are used for variational states.
There is no single resource that reviews the physics-motivated approaches like
backflow, hidden fermions etc., so we mostly look at the corresponding papers.

\subsubsection{Backflow transformations via neural networks}

Following D.\ Luo \& B.K.\ Clark, \emph{Backflow Transformations via Neural
Networks for Quantum Many-Body Wave Functions}.

Recall the definition of backflow introduced by Tocchio et al.\ in ``Role of
backflow correlations for the non-magnetic phase of the $t$-$t'$ Hubbard model'',
Eq.~(4):
\begin{align}
\phi^b_{k\sigma}\big(r_{i,\sigma};\vec x\big)
&=\phi_{k\sigma}\big(r_{i,\sigma}\big)+\sum_j \eta_{ij,\sigma}\,
\phi_{k\sigma}\big(r_{j,\sigma}\big),\\
\eta_{ij,\sigma}&=t\,D_i H_j\,\Theta_{|i-j|,\sigma}, \qquad \text{(L\&C (6))}
\end{align}
with $D_i=n_{i,\uparrow}n_{i,\downarrow}$ and
$H_i=(1-n_{i,\uparrow})(1-n_{i,\downarrow})$ counting doublons and holons.

The $i,j$-dependence of the hopping amplitudes $t_{ij}$ is absorbed in the
variational parameters $\Theta_{|i-j|,\sigma}$; e.g.\ for next-nearest-neighbour
hopping one sets them to zero if $|i-j|\neq 1,2$.

Now we make the ansatz to parametrize the backflow part via a neural network:
\begin{equation}
\phi^b_{k\sigma}\big(r_{i,\sigma};\vec x\big)
=\phi_{k\sigma}\big(r_{i,\sigma}\big)+a^{\mathrm{NN}}_{k i,\sigma}(\vec x) \qquad \text{(L\&C (5))}
\end{equation}
The typical advantage of a NN parametrization is that neural networks are
universal approximators. On top of that, one can show \textbf{explicitly how to
construct the backflow with a NN with two hidden layers}:

\paragraph{Step 1: rewrite $\eta$ as a ReLU.}
Assuming $t\,\Theta_{|i-j|,\sigma}\ge0$,
\begin{equation*}
\eta_{ij,\sigma}=t\,D_iH_j\,\Theta_{|i-j|,\sigma}
=t\,n_{i,\uparrow}n_{i,\downarrow}h_{j,\uparrow}h_{j,\downarrow}\,\Theta_{|i-j|,\sigma},
\end{equation*}
which is nonzero only if
$n_{i,\uparrow}=n_{i,\downarrow}=h_{j,\uparrow}=h_{j,\downarrow}=1$, and hence
\begin{equation*}
\eta_{ij,\sigma}=\mathrm{ReLU}\Big[t\,\Theta_{|i-j|,\sigma}
\big(n_{i,\uparrow}+n_{i,\downarrow}+h_{j,\uparrow}+h_{j,\downarrow}-3\big)\Big].
\end{equation*}

\paragraph{Step 2: input layer.}
The NN shall have the input layer
\begin{align*}
\big(\underbrace{\sigma_1,\dots,\sigma_N}_{\text{spin up}},
\underbrace{\sigma_{N+1},\dots,\sigma_{2N}}_{\text{spin down}}\big),
\qquad \text{with} \\
\sigma_i=2n_{i,\uparrow}-1,\quad \sigma_{i+N}=2n_{i,\downarrow}-1,
\quad i=1,\dots,N .
\end{align*}
Inverting the relation for $\sigma_i$ and inserting it into $\eta_{ij,\sigma}$,
\begin{align*}
&\eta_{ij,\sigma}=\mathrm{ReLU}\Big[\tfrac{1}{2}\,t\,\Theta_{|i-j|,\sigma}
\big(\sigma_i+\sigma_{i+N}-\sigma_j-\sigma_{j+N}-2\big)\Big],  
\end{align*}
with $i,j\in\{1,\dots,N\}$ and $\sigma=\pm1$.

\paragraph{Step 3: first hidden layer.}
Use this to construct $\eta_{ij,\sigma}$ as the following NN (cf.\ \cref{fig:sketch_nn_first_layer}):
\begin{equation}
v_k=\sum_{\ell=1}^{2N}w_{k\ell}\,\sigma_\ell+b_k,
\qquad\text{set}\quad k=\sigma\cdot(i+N\cdot j).
\end{equation}
\begin{equation}
\Rightarrow\quad
v_{\sigma\cdot(i+N\cdot j)}=\sum_{\ell=1}^{2N}w_{\sigma\cdot(i+N\cdot j),\ell}\,
\sigma_\ell+b_{\sigma\cdot(i+N\cdot j)},
\end{equation}
with
\begin{equation}
b_{\sigma\cdot(i+N\cdot j)}=-t\,\Theta_{|i-j|,\sigma}
\qquad\text{and}\qquad
w_{\sigma\cdot(i+N\cdot j),\ell}=\frac{t}{2}\,\Theta_{|i-j|,\sigma}
\begin{cases}
\ \ 1 & \text{if }\ell\ \%\ N=i\\
-1 & \text{if }\ell\ \%\ N=j\\
\ \ 0 & \text{else}
\end{cases}
\end{equation}
such that $\eta_{ij,\sigma}=\mathrm{ReLU}\big(v_{\sigma\cdot(i+N\cdot j)}\big)$.

% ---- sketch of the first hidden layer (requires \usepackage{tikz}) ----------
\begin{figure}[h]
\centering
\begin{tikzpicture}[>=stealth, node distance=1cm]
  \node (in0) at (0,4)    {$1$};
  \node (in1) at (0,3)    {$\sigma_i$};
  \node (in2) at (0,2)    {$\sigma_j$};
  \node (dots) at (0,1.2) {$\vdots$};
  \node (in3) at (0,0.4)  {$\sigma_{i+N}$};
  \node (in4) at (0,-0.6) {$\sigma_{j+N}$};
  \node (out) at (5,1.7)  {$v_k$};
  \foreach \n in {in0,in1,in2,in3,in4}{\fill (\n.east) ++(0.25,0) circle (1.4pt);}
  \draw[->] ($(in0.east)+(0.25,0)$) -- node[above,sloped,font=\small]{$b_{0k}$} (out);
  \draw[->] ($(in1.east)+(0.25,0)$) -- node[above,sloped,font=\small]{$w_{ki}$} (out);
  \draw[->] ($(in2.east)+(0.25,0)$) -- node[above,sloped,font=\small]{$w_{kj}$} (out);
  \draw[->] ($(in3.east)+(0.25,0)$) -- node[below,sloped,font=\small]{$w_{k,i+N}$} (out);
  \draw[->] ($(in4.east)+(0.25,0)$) -- node[below,sloped,font=\small]{$w_{k,j+N}$} (out);
\end{tikzpicture}
\caption{First hidden layer producing $v_k$, with $k=\sigma\cdot(i+N\cdot j)$.}
\label{fig:sketch_nn_first_layer}
\end{figure}

\paragraph{Step 4: second hidden layer.}
Finally, the full backflow is obtained by adding another hidden layer which takes
the $\eta_{ij,\sigma}$ as inputs and weights them by the single-particle orbitals
$\phi_k(r_{j,\sigma})$. The activation function of this layer is simply the
identity. If $t\,\Theta_{|i-j|,\sigma}<0$, the minus sign can be absorbed into
this layer, resulting in
\begin{equation}
a^{\mathrm{NN}}_{k i,\sigma}(\vec r)=\sum_j
\underbrace{\eta_{ij,\sigma}}_{\text{inputs}}\,
\underbrace{s_{ij,\sigma}\,\phi_{k\sigma}(r_{j,\sigma})}_{\text{weights}},
\qquad
s_{ij,\sigma}=\mathrm{sgn}\big(t\,\Theta_{|i-j|,\sigma}\big)=\pm1 .
\end{equation}
\textbf{$\Rightarrow$ the backflow can be represented by a neural network.}

In practice, one simply chooses some NN (RBM, MLP, CNN, \dots) and lets
backpropagation take care of the rest.

\subsection{References}
\begin{itemize}
\item F.\ Schwabl, \emph{Quantum Mechanics}, section 13.1.
\item E.\ Koch, \emph{Mean-Field Theory: Hartree-Fock and BCS}, J\"ulich lecture
      notes (2024), section 1.
\item F.\ Becca \& S.\ Sorella, \emph{Quantum Monte Carlo Approaches for
      Correlated Systems}, sections 1.4-1.5.
\item P.\ Fazekas, \emph{Lecture Notes on Electron Correlation and Magnetism},
      chapter 5 ($t$-$J$ model).
\item L.F.\ Tocchio, F.\ Becca, A.\ Parola \& S.\ Sorella, \emph{Role of backflow
      correlations for the nonmagnetic phase of the $t$-$t'$ Hubbard model},
      Phys.\ Rev.\ B \textbf{78}, 041101(R) (2008).
\item D.\ Luo \& B.K.\ Clark, \emph{Backflow Transformations via Neural Networks
      for Quantum Many-Body Wave Functions}, Phys.\ Rev.\ Lett.\ \textbf{122},
      226401 (2019).
\end{itemize}

\end{document}